\begin{proof}[Proof of \cref{obj:thm:polynomial-net-law-estimator}]
\begin{enumerate}
\item Put
\[
D:=4d_zd_x+d_x,\qquad
b:=4\sqrt{d_zd_x}+\sqrt{d_x}.
\]
By \cref{obj:thm:gap-free-positive-measure-modulus}, choose
\(C_{\mathrm{mod}}>0\) such that, for any two model laws \(P_1,P_2\in\mathcal M\),
\[
W_1(\nu_{P_1},\nu_{P_2})
\le C_{\mathrm{mod}}\,d_S(S(P_1),S(P_2)).
\]
Let \(C_{\mathrm{card}}>0\) and \(C_{\mathrm{work}}>0\) be the finite-library
cardinality and work constants constructed below, and let \(C_0\ge1\) be the
constant supplied by \cref{obj:lem:uniform-summary-concentration}.  Define
\[
C_{\mathrm{tail}}
:=
\max\{C_0,\exp(1),C_{\mathrm{mod}}(2C_0L+1)\},
\qquad
C:=\max\{C_{\mathrm{card}},C_{\mathrm{work}},C_{\mathrm{tail}}\}.
\]
Then \(C>0\), and this single constant dominates the constants used for the
library size, work bound, and probability bound. % lean: finite_net_law_estimator

\item Fix \(n\ge1\), and set
\[
\epsilon_n:=n^{-1/2},\qquad
\delta_n:=\frac{\epsilon_n}{b}.
\]
The library is the deterministic half-open coordinate-grid library in the
\(D\)-coordinate summary box, with side length \(\delta_n\).  Its coordinates
are the entries of the four \(d_z\times d_x\) matrix blocks and the \(d_x\)
entries of the mean block.  For each grid cube meeting the admissible image
\(\mathscr S=S(\mathcal M)\), choose one representative \(S_i\in\mathscr S\)
and order the representatives by the lexicographic order of the cube lower
corners.  The coordinate bounds following from
\cref{obj:prop:observed-vmw-margin-inclusion} place every admissible summary
coordinate in \([-L,L]\).  If \(q\in\mathscr S\) and \(S_i\) lie in the same
cube, then every scalar coordinate differs by at most \(\delta_n\), so
\[
\|M_t(S_i)-M_t(q)\|_{\mathrm{op}}
\le \sqrt{d_zd_x}\,\delta_n,\qquad
\|N_t(S_i)-N_t(q)\|_{\mathrm{op}}
\le \sqrt{d_zd_x}\,\delta_n
\]
for \(t\in\{0,1\}\), and
\[
\|m_X(S_i)-m_X(q)\|_2\le \sqrt{d_x}\,\delta_n .
\]
Using \(d_S\) from \cref{obj:def:observable-summary},
\[
d_S(S_i,q)
\le
\bigl(4\sqrt{d_zd_x}+\sqrt{d_x}\bigr)\delta_n
=\epsilon_n .
\]
Thus the stored representatives form an \(n^{-1/2}\)-net of \(\mathscr S\);
when \(\mathscr S=\varnothing\), the index set is empty and the program uses
the fallback law in \cref{obj:def:finite-net-program}. % lean: netLibrary_nonempty

\item Each coordinate lower corner has at most
\[
\left\lceil \frac{2L}{\delta_n}\right\rceil+1
\]
possible values.  Since \(\delta_n=n^{-1/2}/b\) and \(n^{1/2}\ge1\),
\[
\left\lceil \frac{2L}{\delta_n}\right\rceil+1
\le (2Lb+2)n^{1/2}.
\]
The lower-corner code is injective, hence
\[
|A|
\le \bigl((2Lb+2)n^{1/2}\bigr)^D
=(2Lb+2)^D n^{D/2}
\le C_{\mathrm{card}} n^{(4d_zd_x+d_x)/2}.
\]
After enlarging \(C\), the asserted cardinality bound follows. % lean: netLibrary_card_polynomial_bound

\item Fix a uniform provider of the exact-real primitives.  The estimator is the
law component of \(\operatorname{NetProgram}(A,\omega)\) in
\cref{obj:def:finite-net-program}.  On a nonempty library, enumerate the finite
index set as \(\iota_0,\ldots,\iota_{m-1}\) and update the current index by
choosing the smaller pair
\[
\bigl(d_S(S_i,s),\operatorname{rank}_{\mathrm{lex}}(i)\bigr)
\]
in lexicographic order.  For each fixed representative,
\(s\mapsto d_S(S_i,s)\) is continuous by
\cref{obj:lem:summary-metric-continuity}; therefore every update is Borel,
being defined by finitely many weak and strict Borel inequalities.  The empty
selection is constant.  Since \(\widehat S_n\) is Borel by
\cref{obj:def:empirical-summary}, and the selected finite index is mapped to a
stored spectral law, \(\widehat\nu_n^{\mathrm{net}}\) is Borel measurable. % lean: netLawEstimator_measurable

\item The exact-real trace is assembled from the empirical-summary trace, the
finite-library search trace, and the selected spectral run:
\[
\mathfrak T(\omega)
=
\mathfrak T_{\mathrm{emp}}(n)
\,\Vert\,
\mathfrak T_{\mathrm{search}}(A)
\,\Vert\,
\begin{cases}
\mathfrak T_{\mathrm{spec}}(S_i),
&\text{if the search selects representative }i,\\
\varnothing,
&\text{if the search has empty selection.}
\end{cases}
\]
Thus the returned law and the displayed trace come from the same execution. % lean: netExactRealProgram_trace_eq

\item The empirical-summary trace has \(nD\) arithmetic operations.  Each
library comparison uses \(D\) scalar arithmetic operations, four
fixed-dimensional singular-value or operator-norm operations, and one
comparison.  The selected spectral tail has three singular-value operations,
one root-isolation operation, and one arithmetic operation.  Therefore, for
every sample \(\omega\),
\[
\operatorname{ops}(\omega)
\le
nD+|A|(D+5)+5 .
\]
Combining this with the cardinality bound gives
\[
\operatorname{ops}(\omega)
\le
C_{\mathrm{work}}\{n+n^{D/2}\}
\le
C\{n+n^{(4d_zd_x+d_x)/2}\}.
\] % lean: netOperationCount_polynomial_bound

\item If the program selects \(i\) at sample \(\omega\), put
\[
s_\omega:=\widehat S_n(\omega).
\]
The exhaustive fold preserves domination, where \(i\) dominates \(j\) at
\(s_\omega\) when
\[
d_S(S_i,s_\omega)\le d_S(S_j,s_\omega)
\]
and equality implies
\[
\operatorname{rank}_{\mathrm{lex}}(i)
\le
\operatorname{rank}_{\mathrm{lex}}(j).
\]
Induction over the scanned list yields, for every library index \(j\),
\[
d_S(S_i,\widehat S_n(\omega))
\le d_S(S_j,\widehat S_n(\omega)),
\]
and
\[
d_S(S_i,\widehat S_n(\omega))
=
d_S(S_j,\widehat S_n(\omega))
\Longrightarrow
\operatorname{rank}_{\mathrm{lex}}(i)
\le
\operatorname{rank}_{\mathrm{lex}}(j).
\] % lean: nearestLibraryIndex_spec

\item Let \(\mathcal R_i\) be the exact spectral run at the selected
representative \(S_i\).  By the definition of the program output,
\begin{equation}
\widehat\nu_n^{\mathrm{net}}(\omega)=\operatorname{law}(\mathcal R_i).
\label{eq:finite-net-selected-run-law}
\end{equation} % lean: netLawEstimator_eq_of_selected

\item The same spectral execution supplies a basis \(V_i\), its signal spans,
and an eigenvalue list \((\lambda_{i,r})_{r<k}\).  Its output fields give, for
each \(t\in\{0,1\}\),
\[
x\longmapsto M_t(S_i)V_i x
\]
injective, the exact threshold characterization
\[
\frac{\pi_0\sigma_0^2}{2}
\le
\sigma_j\!\left(\operatorname{stack}(M_0(S_i),M_1(S_i))\right)
\Longleftrightarrow j<k,
\]
and completeness of the returned eigenvalue list:
\[
z\in\operatorname{spec}(\Delta_i)
\Longrightarrow
\exists r<k\ \text{such that }z=\lambda_{i,r},
\]
where \(\Delta_i\) is the compressed operator formed from \(S_i\), \(V_i\), and
the run's spans. % lean: ExactRealSpectralRun.output

\item Fix any stored representative \(S_i\).  Since \(S_i\in\mathscr S\), choose
\(Q_i\in\mathcal M\) with
\[
S(Q_i)=S_i.
\]
Let \(p_{i,u}:=Q_i(U=u)\), write \(p_i:=(p_{i,u})_{u=1}^k\), and set
\[
\tau_{i,u}:=\mu_{1u}(Q_i)-\mu_{0u}(Q_i).
\]
The latent-arm positivity condition in \cref{obj:ass:latent-arm-positivity}
gives \(p_{i,u}>0\) for every \(u\).  Let \(\mathcal R_i\) be the exact
spectral run at \(S_i\), with signal basis \(V_i\), spans, anchors \(a_i,c_i\),
returned eigenvalues \((\lambda_{i,r})_{r<k}\), and compressed operator
\[
\Delta_i
:=(M_1(S_i)V_i)^\dagger N_1(S_i)V_i
 -(M_0(S_i)V_i)^\dagger N_0(S_i)V_i .
\]
Set
\[
\Omega_i:=B(Q_i)^\top V_i\in\mathbb R^{k\times k},
\qquad
\Gamma_{i,t}
:=\operatorname{diag}\{Q_i(U=u\mid T=t):1\le u\le k\},
\]
and
\[
\Phi_{i,t}:=A_t(Q_i)\Gamma_{i,t}\in\mathbb R^{d_z\times k}.
\]
\Cref{obj:lem:observed-proxy-moment-factorization,obj:lem:observed-outcome-proxy-moment-factorization}
give, for \(t\in\{0,1\}\),
\[
M_t(S_i)V_i=\Phi_{i,t}\Omega_i,
\qquad
N_t(S_i)V_i
=\Phi_{i,t}
\operatorname{diag}(\mu_{t1}(Q_i),\ldots,\mu_{tk}(Q_i))\Omega_i .
\]
The run's armwise full-rank certificate makes
\(x\mapsto M_t(S_i)V_ix\) injective for each \(t\).  Taking \(t=0\), the first
factorization shows that \(\Omega_i\) is injective and hence invertible.  The
same factorization also gives injectivity of \(\Phi_{i,0}\): if
\(\Phi_{i,0}v=\Phi_{i,0}v'\) for \(v,v'\in\mathbb R^k\), then
\[
M_0(S_i)V_i\Omega_i^{-1}v
=
\Phi_{i,0}v
=
\Phi_{i,0}v'
=
M_0(S_i)V_i\Omega_i^{-1}v',
\]
so injectivity of \(M_0(S_i)V_i\) gives
\(\Omega_i^{-1}v=\Omega_i^{-1}v'\), and hence \(v=v'\).  The proxy
factorization before postmultiplication by \(V_i\) is
\[
M_0(S_i)=\Phi_{i,0}B(Q_i)^\top .
\]
Since \(\Phi_{i,0}\) is injective, its Moore--Penrose inverse is a left inverse
on its range, and therefore
\[
\Phi_{i,0}^\dagger M_0(S_i)=B(Q_i)^\top .
\]
Transposing gives
\[
B(Q_i)=M_0(S_i)^\top(\Phi_{i,0}^\dagger)^\top .
\]
By construction of the run's span certificate,
\[
\operatorname{range}(V_i)
=
\operatorname{range}(M_0(S_i)^\top)+\operatorname{range}(M_1(S_i)^\top)
\subseteq\mathbb R^{d_x},
\]
and consequently
\[
\operatorname{range}(B(Q_i))
\subseteq
\operatorname{range}(M_0(S_i)^\top)
\subseteq
\operatorname{range}(V_i).
\]
The second factorization can therefore be rewritten as
\[
N_t(S_i)V_i
=
M_t(S_i)V_i
\{\Omega_i^{-1}
\operatorname{diag}(\mu_{t1}(Q_i),\ldots,\mu_{tk}(Q_i))\Omega_i\}.
\]
Since the Moore--Penrose inverse is a left inverse on the range of an injective
linear map,
\[
(M_t(S_i)V_i)^\dagger N_t(S_i)V_i
=
\Omega_i^{-1}
\operatorname{diag}(\mu_{t1}(Q_i),\ldots,\mu_{tk}(Q_i))\Omega_i .
\]
Subtracting the two treatment-arm identities gives
\[
\Delta_i
=
\Omega_i^{-1}
\operatorname{diag}(\tau_{i,1},\ldots,\tau_{i,k})
\Omega_i .
\]
Finally, the law of total expectation gives
\[
m_X(S_i)=B(Q_i)p_i .
\]
The column-space inclusion just proved implies that the orthogonal projector
\(V_iV_i^\top\) fixes \(B(Q_i)\).  Since
\(\Omega_i^\top=V_i^\top B(Q_i)\),
\[
B(Q_i)=V_iV_i^\top B(Q_i)=V_i\Omega_i^\top .
\]
Thus
\[
a_i=V_i^\top m_X(S_i)=\Omega_i^\top p_i .
\]
The target-proxy anchor condition gives \(B(Q_i)^\top e_1=\mathbf 1_k\), while
\(c_i=V_i^\top e_1\) and \(B(Q_i)^\top V_iV_i^\top=B(Q_i)^\top\), so
\[
\Omega_i c_i=\mathbf 1_k .
\] % lean: modelCompressedCoordinates_exists
For the polynomial projector attached to a returned slot \(r\),
\[
\mathcal P_{i,r}
:=
\prod_{\lambda_{i,\ell}\ne\lambda_{i,r}}
(\lambda_{i,r}-\lambda_{i,\ell})^{-1}
(\Delta_i-\lambda_{i,\ell}I_k),
\]
the diagonalization gives
\[
\mathcal P_{i,r}
=
\begin{cases}
\Omega_i^{-1}
\operatorname{diag}\{\mathbf 1(\tau_{i,u}=\lambda_{i,r})\}_{u=1}^k
\Omega_i,
&\lambda_{i,r}\in\{\tau_{i,u}:1\le u\le k\},\\
0,&\lambda_{i,r}\notin\{\tau_{i,u}:1\le u\le k\}.
\end{cases}
\]
Hence the returned slot weight
\[
w_{i,r}:=a_i^\top\mathcal P_{i,r}c_i
\]
satisfies
\[
w_{i,r}
=
\sum_{u:\tau_{i,u}=\lambda_{i,r}}p_{i,u}
\]
for a spectral value and \(w_{i,r}=0\) for a listed value outside the spectrum.
The output validity of \(\mathcal R_i\) gives \(w_{i,r}\ge0\) and
\(\sum_{r<k}w_{i,r}=1\).  Completeness gives at least one returned slot for
each latent-effect value, and the positivity of the corresponding cluster mass
\[
m_i(\lambda):=\sum_{u:\tau_{i,u}=\lambda}p_{i,u}
\]
forces exactly one positive returned slot for each distinct latent-effect
value; otherwise the total \(\sum_{r<k}w_{i,r}\) would exceed
\(\sum_{\lambda}m_i(\lambda)=1\).  Therefore, for every bounded test function
\(\varphi:\mathbb R\to\mathbb R\),
\[
\sum_{r<k}w_{i,r}\varphi(\lambda_{i,r})
=
\sum_{u=1}^k p_{i,u}\varphi(\tau_{i,u}).
\]
Thus the induced atomic measures agree:
\[
\sum_{r<k}w_{i,r}\delta_{\lambda_{i,r}}
=
\sum_{u=1}^kp_{i,u}\delta_{\tau_{i,u}}
=
\nu_{Q_i},
\]
with coincident effects aggregated as in \cref{obj:def:quotient-law}.  The exact
spectral run at \(S_i\) returns \(\nu_{Q_i}\):
\begin{equation}
\operatorname{law}(\mathcal R_i)=\nu_{Q_i}.
\label{eq:finite-net-run-quotient}
\end{equation} % lean: exactRealSpectralRun_eq_quotient

\item The validity certificate of the same result-bearing spectral execution
states
\[
w_{i,r}\ge0,\qquad
\sum_{r<k}w_{i,r}=1,\qquad
|\lambda_{i,r}|\le L_\tau
\quad(r<k).
\]
Consequently the returned weights form a probability vector and every returned atom lies in \([-L_\tau,L_\tau]\), so the represented finite atomic measure is a probability law supported on that interval. % lean: exactRealSpectralRun_output_valid

\item For a stored representative \(S_i=S(Q_i)\), put
\[
\mathcal G_i:=\operatorname{stack}(M_0(S_i),M_1(S_i)),
\qquad
s_0:=\pi_0\sigma_0^2 .
\]
\Cref{obj:lem:model-compressed-spectral-certificate} supplies the rank identity,
and \cref{obj:lem:stacked-proxy-moment-signal-margin} supplies the
singular-value margin:
\[
\operatorname{rank}(\mathcal G_i)=k,
\qquad
s_0\le \sigma_{k-1}(\mathcal G_i).
\]
If \(H\in\mathbb R^{2d_z\times d_x}\) satisfies
\[
\|H\|_{\mathrm{op}}<s_0/2,
\]
Weyl's singular-value inequality gives, for every \(r\),
\[
|\sigma_r(\mathcal G_i+H)-\sigma_r(\mathcal G_i)|
\le \|H\|_{\mathrm{op}}.
\]
Thus
\[
\sigma_r(\mathcal G_i+H)>s_0/2\quad(r<k),
\]
while \(\operatorname{rank}(\mathcal G_i)=k\) gives
\(\sigma_r(\mathcal G_i)=0\) for \(r\ge k\), and hence
\[
\sigma_r(\mathcal G_i+H)<s_0/2\quad(r\ge k).
\]
Consequently the perturbed stacked moment satisfies the two threshold clauses
\[
\frac{s_0}{2}\le \sigma_r(\mathcal G_i+H)\quad(r<k),
\qquad
\sigma_r(\mathcal G_i+H)<\frac{s_0}{2}\quad(r\ge k).
\]
Equivalently, exactly \(k\) singular values of \(\mathcal G_i+H\) meet the
threshold \(s_0/2=\pi_0\sigma_0^2/2\). % lean: modelSummary_thresholdRecovers_of_perturbation

\item Fix \(P\in\mathcal M\) and a sample \(\omega\).  If \(P\in\mathcal M\),
then \(\mathscr S\) is nonempty, so the library has a selected nearest index
\(i\).  By the covering property, choose \(j\) such that
\[
d_S(S_j,S(P))\le n^{-1/2}.
\]
Nearest-library optimality and the triangle inequality yield
\[
\begin{aligned}
d_S(S_i,S(P))
&\le d_S(S_i,\widehat S_n(\omega))
   +d_S(\widehat S_n(\omega),S(P))\\
&\le d_S(S_j,\widehat S_n(\omega))
   +d_S(\widehat S_n(\omega),S(P))\\
&\le d_S(S_j,S(P))
   +2d_S(\widehat S_n(\omega),S(P))\\
&\le n^{-1/2}+2d_S(\widehat S_n(\omega),S(P)).
\end{aligned}
\]
By \cref{eq:finite-net-selected-run-law,eq:finite-net-run-quotient},
\[
\widehat\nu_n^{\mathrm{net}}(\omega)=\nu_{Q_i},
\qquad
S(Q_i)=S_i.
\]
Applying the modulus from \cref{obj:thm:gap-free-positive-measure-modulus}
gives
\[
W_1\!\left(\widehat\nu_n^{\mathrm{net}}(\omega),\nu_P\right)
\le
C_{\mathrm{mod}}
\left\{
2d_S(\widehat S_n(\omega),S(P))+n^{-1/2}
\right\}.
\] % lean: netEstimator_wass1_le

\item Fix \(0<\eta<1/2\).  Since \(C\ge C_0\), \(C\ge\exp(1)\), and
\(\eta\le1\),
\[
1\le \log(C/\eta),
\qquad
\log(C_0/\eta)\le \log(C/\eta).
\]
Therefore
\[
n^{-1/2}\le
\sqrt{\frac{\log(C/\eta)}{n}},
\qquad
\sqrt{\frac{\log(C_0/\eta)}{n}}
\le
\sqrt{\frac{\log(C/\eta)}{n}}.
\]
On the event
\[
d_S(\widehat S_n,S(P))
\le
C_0L\sqrt{\frac{\log(C_0/\eta)}{n}},
\]
the deterministic inequality gives
\[
W_1(\widehat\nu_n^{\mathrm{net}},\nu_P)
\le
C_{\mathrm{mod}}(2C_0L+1)
\sqrt{\frac{\log(C/\eta)}{n}}
\le
C\sqrt{\frac{\log(C/\eta)}{n}}.
\]
Thus the bad event in the theorem is contained in
\[
\left\{
d_S(\widehat S_n,S(P))
>
C_0L\sqrt{\frac{\log(C_0/\eta)}{n}}
\right\}.
\]
By \cref{obj:lem:uniform-summary-concentration}, its \(Q_P^{(n)}\)-probability
is at most \(\eta\). % lean: finite_net_law_estimator
\end{enumerate}
\end{proof}
